% Français 9115, batch 26 — Gallica views 501–520.
% Pass: Opus 5.5 (claude-opus-5-5), 2026-09-22 — first pass, unchecked against the views by a human.
%
% What the twenty views hold. One piece, begun before the batch and going
% on after it, in French, in her regular fair hand, corrected in the usual
% way (words struck on the line, replacements written above it in a finer
% pen, two marginal notes): pages (91) to (110) of the fair copy of her
% prize memoir on vibrating elastic surfaces — the memoir whose pages
% (1)–(34) were read in batches 20–22 (v. 398–440). Every view of the
% batch carries text: the leaves are written on both sides, recto and
% verso, each mounted on a guard sheet (on many views the blank top or
% foot of the guard occupies a quarter of the frame). No blank verso, no
% repeat, no microfilm target, no printed matter, no figure drawn.
%
%   v. 501  (91) end of the comparison of the square plates' sounds with
%           the theory: the sounds of the figures with complete periods
%           are about one tone graver than the theory gives (« 0|4 fig.
%           72–73 » sol♯3 against « 4|3 fig. 80 » fa♯5, two octaves
%           expected, 25/4 against 25).
%   v. 502  (92) the difference is a minimum, approached as the nodal lines
%           grow in number: figs 63, 68, 66a, 69, 74, 79b, 86, 75; « V
%           table P. 152 ».
%   v. 503  (93) the same on rectangular plates; the differences rarely
%           exceed a third, but their constancy forbids ascribing them to
%           the physical causes of experiment.
%   v. 504  (94) her conjecture: « quart de période », « demie période »;
%           the difference greatest for quarter-periods, least with
%           complete periods. The page ends two thirds down.
%   v. 505  (95) the quarter- and half-periods as the cause; only plates
%           supported all round would agree exactly, and Chladni never
%           tried them; a lamina free by ddz/dx²=0, d³z/dx³=0 would have
%           decided it.
%   v. 506  (96) « une simple conjecture »; the string and the open flute
%           are no help, the theory of wind instruments not being rigorous;
%           « M^r De la Place », « M^r Poisson » and the velocity of sound.
%   v. 507  (97) the heat coefficient and the half-periods would be
%           confounded; « soumettre au jugement dela classe »; the
%           difference barely exceeds those Chladni neglected in his
%           « table p. 152 ». Ends « Je vais m'occuper àprésent des plaques
%           rectangulaires differentes du quarré. », the rest of the page
%           blank.
%   v. 508  (98) new section, underlined: « §.5 Comparaison dela théorie à
%           l'experience dans le cas des plaques rectangulaires differentes
%           du quarré »; N.o 17, integral (C) N.o 3 (« V. N.os 10 11 et 12
%           §.4 »); ω the same in x and y, so B = 0 or scriptA = 0 in most
%           cases; the plate 5:7, numbers 25/4 and 49/4.
%   v. 509  (99) 3 lines one way and 4 the other give the same sound;
%           Chladni's « figures intermédiaires », traité d'acoustique
%           N.o 136 and N.o 121; the plate 3:5.
%   v. 510  (100) the plate 3:1, cases 5|0 and 0|2, fig. 77 (« T.
%           d'acoustique N 122 »); a marginal note replacing a struck
%           passage on the figures of « pl. 6 » and fig. 170; a struck
%           paragraph announcing fig. 163b.
%   v. 511  (101) f(x), f(y) at the origin of the form 2(1+e^ω) or
%           2(1−e^ω); the supposition scriptA = −B gives no parallel and
%           diagonal lines except on the square.
%   v. 512  (102) fig. 175b (2|0 and 0|4); for an odd and an even number,
%           scriptA(1−e^ω) = −(1+e^ω)B; fig. 163b traced from the point a
%           to the point i.
%   v. 513  (103) the four branches round a point m (« N.o 11 »); the 12
%           « points de repos invariables »; the other figures 163,
%           « flexions » and half-flexions.
%   v. 514  (104) the general rule: m lines for scriptA = 0, n for B = 0,
%           so m sinuous lines of n/2 flexions or n of m/2; N.o 18
%           « surfaces periodiques », integrals N.os 4 and 5.
%   v. 515  (105) z = … sin(π.tmx/tA) sin(π.sny/tA) for sides s:t; the
%           sound as (t²m²+s²n²)/t²; Chladni did not publish all his
%           rectangular figures.
%   v. 516  (106) sums of particular integrals, admissible only when every
%           term gives the same sound.
%   v. 517  (107) invariable points and true invariable lines of rest;
%           « libre » again has neither of its two meanings.
%   v. 518  (108) s = 6, t = 5: 4(109)/25 against 4(106)/25, an interval
%           109/106, « guere plus grand qu'un comma », below the errors of
%           experiment.
%   v. 519  (109) the central line of rest (« V. N.o 114 fig. 157 »); a
%           marginal note on the rotary vibrations; the plate 5:4, figures
%           5|0 and 1|4 (« V. N.o 115 »).
%   v. 520  (110) figs 158; s = 5, t = 4, m = 1, n = 4 gives 26; Euler's
%           number for 5|0, (81/4)(25/16); the sounds :: 2025 : 1664,
%           about 5:4, two tones; « figures compliquées par leur influence
%           simultanée ». Stops mid-sentence.
%
% Seams. View 501 opens mid-sentence: « et dans celui des y, on prend pour
% terme de comparaison… »; the thumbnail of view 500 (not transcribed)
% ends « …des periodes complettes dans le sens des x », so the sentence
% runs across 500/501. View 520 stops on « 5 lignes nodales droites
% parallelles entr'elles » followed by a struck « et »; view 521, page
% (111), ink 256, opens with a struck line and then « et au plus grand côté
% dela plaque », so the text runs across 520/521. Inside the batch the text
% is continuous from page to page; new paragraphs begin at v. 504 (a page
% ending two thirds down), v. 508 (the new section, after v. 507 ends on a
% blank third) and v. 516. References outside the batch: « l'intégrale
% (C) » and « (I) », « N.o 3 », « N.os 10 11 et 12 §.4 », « N.o 11 »,
% « N.os 4 et 5 », « la table » and « table p. 152 », and Chladni's traité
% d'acoustique by paragraph (N.os 114, 115, 121, 122, 136) and plate
% (« pl. 6 »).
%
% The numbers on the leaves, and why no \folio{} is written. (1) Her
% pagination, in parentheses at the head, centre: (91) at v. 501, then
% (92) to (110) at v. 502–520, one per view, without a gap and none struck.
% (2) The ink series at the top right, in the larger rounder hand: 246
% (v. 501), 247 (503), 248 (505), 249 (507), 250 (509), 251 (511), 252
% (513), 253 (515), 254 (517), 255 (519) — one per leaf, on the odd views,
% which her pagination makes odd pages; none on the even views, which are
% the backs of the same leaves. It goes on to 256 at v. 521. As in every
% batch of this volume it is penned, not pencilled, and no \folio{} is
% written. (3) Her paragraph numbers: §.5 and N.o 17 (v. 508), N.o 18
% (v. 514). (4) Chladni's figure numbers (0|4, 4|3, 5|0, 0|2, 2|0, 1|4 …)
% and numbered figures (63–86, 157, 158, 163, 163b, 170, 175b), which are
% his. Every number above was read on its view; none is inferred. None of
% these leaves can be Laubenbacher and Pengelley's « Manuscript C »
% (ff. 348r–349r): the subject is vibrating plates, and the ink series
% here runs 246–255.
%
% Notation. Hers, kept. z the ordinate, ω the angle of Euler's cases,
% m and n the numbers of nodal lines, s:t the sides of the plate; letters
% she underlines in the prose (z, x, y, m, n, ω, a) keep \underline{} in
% maths; words she underlines are set in \emph{}, and so are the
% underlined letter of « 163b », « 175b » and a struck « 3 » in the prose
% (text-mode \underline fails the TEI check). The coefficient paired
% with B is the looped cursive capital of batch 22, set $\mathcal{A}$.
% Chladni's figure signs « 0|4 » are set $0|4$. « sin » is set \text{sin};
% the points « z = …… » are hers, a factor left blank. Proportions are
% written « :: 5:7 » and kept. Notes of music are written « sol♯3 »,
% « fa♯5 », « ut.6 », set with $^{\sharp}$. Plurals cancelled by the small
% « x » on the final s (hand.md, v. 314) are given in the singular with a
% note. Where the leaf and the calculation disagree the leaf stands, with a
% note: v. 518 (« 96 » where 36 is required), v. 520 (a numerator
% overwritten, a struck intermediate fraction). Spelling is hers and is not
% flagged: « complettes », « proportionel », « guères », « quarré »,
% « parallelles », « differens », « mouvemens », « instrumens », « tems »,
% « satisfesantes », « fesant », « atribue », « embarassé », « devroit »,
% « étoit », « paroît », « connoître », « aux premiers » for aux premières,
% « proportionnés », « transversales », and the run-together words
% « dela », « delasimple », « desorte », « cequ'il », « ceque »,
% « aumoins », « àprésent », « lamême », « laformation », « delignes »,
% « àlaquelle », « àlaquestion », « toutdesuite », « àlafois », « quelquesorte »,
% « c'est adire ».
%
% What could not be read, and what is offered with doubt. \ill{} stands
% thirteen times: a blackened group after « (C) et » (v. 501); the middle
% of a struck group after « (I), » (v. 502); a struck word after « il »
% and a blackened group after « desorte, » (v. 504); the end of a struck
% line after « appuyer » (v. 505); a struck word after « excède » (v. 507);
% a struck group after « de » in the marginal note (v. 510); a struck word
% before « nombre » in the marginal note and a letter after « qu » (v. 519);
% the first factor of a numerator, the numerator of a struck fraction, a
% struck word before « doivent égaux » and a struck group after « plus
% graves » (v. 520). \uncertain{} stands for « au » (v. 503), « m'avoit »
% (v. 505), « ont » (v. 510), « et » (v. 511), « pourra » (v. 517),
% « Au reste », « 4|1 » and « sait » in the marginal note and « cit »
% (v. 519).
%
% Nothing here was checked against any published edition of Chladni, of
% Euler, of Laplace or Poisson, or of Sophie Germain's papers.
\documentclass[11pt,a4paper]{article}
% The preamble shared by every transcription of the mathematicians' manuscripts in Gallica.
%
% It rests on one idea: **uncertainty compounds, it does not resolve.** A
% transcription that smooths over a doubtful word has destroyed the only thing
% it was made for — a reader must be able to tell, at every point, what was on
% the page and what was supplied. The apparatus macros below are therefore the
% critical apparatus, not decoration.
%
% The same commands are recognised by `scripts/render.mjs`, which produces the
% site's reading view, and by `scripts/tei.mjs`. A macro added here must be
% added there too, or rendering fails loudly — which is the wanted behaviour.
%
% Comments here are English, like the rest of the repository. The documents
% this preamble sets are French, and that is deliberate: see the skills.

\ProvidesPackage{germain}[2026/09/15 Transcriptions of mathematicians' manuscripts in Gallica]

\usepackage{amsmath,amssymb,amsthm}
\usepackage{mathrsfs}
% Kept from the parent preamble so the renderer's diagram support stays
% available; these manuscripts draw no commutative diagrams, and nothing here uses it.
\usepackage{tikz-cd}
\usepackage[english,french]{babel}
\usepackage{fontspec}

% --- Greek ------------------------------------------------------------------
% The text face stays fontspec's default, Latin Modern, which has no Greek:
% XeTeX drops every glyph it lacks with a line in the log and nothing on the
% page, and the Greek words the manuscripts quote vanished from the PDFs. So
% the Greek and Greek Extended blocks get a XeTeX character class of their own,
% and entering or leaving a run of it switches the family to CMU Serif — the
% same Computer Modern design, with polytonic Greek — and back. Only the
% family changes: a Greek word in \textit or \textbf keeps its shape and series.
% The transitions to and from every other class carry the tokens that class
% has towards an ordinary letter, so French spacing before « ; » or inside
% « » still holds after a Greek word. No grouping, so no brace or \endgroup
% can come out unbalanced; maths is left alone, since XeTeX inserts nothing
% there. Kept identical in `exercices.sty`.
\makeatletter
\newfontfamily\germainGreekfont{cmun}[
  Extension=.otf, NFSSFamily=germaingreek,
  UprightFont=*rm, ItalicFont=*ti, SlantedFont=*sl,
  BoldFont=*bx, BoldItalicFont=*bi, BoldSlantedFont=*bl]
\newXeTeXintercharclass\germain@greek
\def\germain@greekrange#1#2{%
  \@tempcnta=#1\relax
  \loop
    \XeTeXcharclass\@tempcnta=\germain@greek
    \ifnum\@tempcnta<#2\advance\@tempcnta\@ne\repeat}
\germain@greekrange{"0370}{"03FF}
\germain@greekrange{"1F00}{"1FFF}
\def\germain@greekprev{\rmdefault}
\def\germain@greekin{%
  \edef\germain@greekprev{\f@family}\fontfamily{germaingreek}\selectfont}
\def\germain@greekout{\fontfamily{\germain@greekprev}\selectfont}
\def\germain@greekjoin{%
  \XeTeXinterchartokenstate=\@ne
  \@tempcnta=\z@
  \loop
    \ifnum\@tempcnta=\germain@greek\else
      \XeTeXinterchartoks\germain@greek\@tempcnta=\expandafter{%
        \expandafter\germain@greekout\the\XeTeXinterchartoks\z@\@tempcnta}%
      \XeTeXinterchartoks\@tempcnta\germain@greek=\expandafter{%
        \the\XeTeXinterchartoks\@tempcnta\z@\germain@greekin}%
    \fi
    \ifnum\@tempcnta<4095\advance\@tempcnta\@ne\repeat}
% babel-french sets its own classes, and saves and restores the interchar
% state, each time a language is selected: join again after it does.
\addto\extrasfrench{\germain@greekjoin}
\addto\extrasenglish{\germain@greekjoin}
\AtBeginDocument{\germain@greekjoin}
\makeatother

\usepackage{geometry}
\geometry{a4paper,margin=25mm,marginparwidth=28mm}
\usepackage{marginnote}
\usepackage{xcolor}
\usepackage[normalem]{ulem}
\setlength{\emergencystretch}{3em}
\usepackage{hyperref}
\hypersetup{colorlinks=true,linkcolor=black,urlcolor=black,pdfborder={0 0 0}}

% --- Batch metadata ---------------------------------------------------------
% Taken as they stand by the rendering script, which shows them at the head of
% the reading view. They fix what the file claims to cover. The macro names are
% the parent project's, so that the scripts stay shared; \folder here means the
% volume's slug — `fr-9115` — and \pages counts Gallica views.

\newcommand{\folder}[1]{\gdef\grothFolder{#1}}
\newcommand{\batch}[1]{\gdef\grothBatch{#1}}
\newcommand{\pages}[2]{\gdef\grothFirst{#1}\gdef\grothLast{#2}}
\newcommand{\foldertitle}[1]{\gdef\grothFoldertitle{#1}}

% The holder's own shelfmark, printed as they write it: « Français 9115 ».
\newcommand{\shelfmark}[1]{\gdef\grothShelfmark{#1}}

% The dating, copied verbatim from the holder's catalogue — never inferred
% here. « XIXe siècle » is the BnF's; a pass that dates a leaf from its content
% says so in a \note{}, as its own inference.
\newcommand{\dating}[1]{\gdef\grothDating{#1}}

% The status notice, printed in the title block and repeated as a diagonal
% watermark on every page of the PDF. The manuscripts are in the public
% domain; what this line declares is not a right but a quality — an
% unverified first pass by a machine. The skills write the line; a file
% without it gets no watermark, so leaving it out is visible in the output.
\newcommand{\watermark}[1]{\gdef\grothWatermark{#1}}

\gdef\grothFolder{?}\gdef\grothBatch{}\gdef\grothFirst{?}\gdef\grothLast{?}%
\gdef\grothFoldertitle{}\gdef\grothShelfmark{}\gdef\grothDating{}\gdef\grothWatermark{}

% --- The résumé -------------------------------------------------------------
\newenvironment{resume}
  {\small\setlength{\parskip}{0.45em}}
  {\par\medskip\hrule\medskip}

% --- Keywords ---------------------------------------------------------------
% In English, closing the modernised reading's résumé: the modern vocabulary
% under which to search for what the leaves construct. The manifest extracts
% them to tag the volume — this line is the single source of the tags.
\newcommand{\keywords}[1]{\par\smallskip\noindent
  {\footnotesize\textsc{Keywords} --- \textit{#1}}\par}

% --- The critical apparatus -------------------------------------------------

% The Gallica view number, in the margin. This is the anchor that lets a
% disagreement over a formula be settled by looking at *one* image rather than
% a volume of 750. The site uses it to turn the facsimile as the transcription
% is scrolled. A view is one scanned image: usually one side of a leaf.
\newcommand{\page}[1]{%
  \marginnote{\footnotesize\color{gray!70}\textsf{v.\,#1}}%
  \label{page:#1}\ignorespaces}

% The BnF's pencilled foliation, where the leaf carries it and a pass can read
% it: « 198r ». This is what the literature cites — « ff. 198r–208v » — and
% the one line that turns a citation into a view. Recorded, never inferred:
% a folio a pass cannot read is not written.
\newcommand{\folio}[1]{%
  \marginnote{\footnotesize\color{gray!50}\textsf{f.\,#1}}\ignorespaces}

% The modernised reading's coarser counterpart to \page{}: which views a
% section is drawn from.
\newcommand{\pagerange}[2]{%
  \marginnote{\footnotesize\color{gray!70}\textsf{v.\,#1--#2}}%
  \ignorespaces}

% Illegible: never guessed.
\newcommand{\ill}{\textcolor{red!60!black}{[\,\ldots\,]}}

% A doubtful reading: the word is given, and flagged as uncertain.
\newcommand{\uncertain}[1]{\uline{#1}}

% An editorial addition — a missing letter, an implied word.
\newcommand{\add}[1]{\textcolor{blue!50!black}{[#1]}}

% Struck out by the writer. What was crossed out is part of the document.
\newcommand{\struck}[1]{\sout{#1}}

% The transcriber's note, on the state of the page or the mathematics.
\newcommand{\note}[1]{\par\smallskip{\small\color{gray!60!black}\textit{[#1]}}\par\smallskip}

% A marginal note of hers, distinct from the body of the text.
\newcommand{\marginal}[1]{\par\smallskip\noindent\hspace{1em}%
  {\small\color{gray!60!black}#1}\par\smallskip}

% --- Title ------------------------------------------------------------------
% The title block is French, like the documents themselves.

\AddToHook{shipout/background}{%
  \ifx\grothWatermark\empty\else
    \begin{tikzpicture}[remember picture, overlay]
      \node[rotate=52, scale=3.4, align=center, text=gray!18,
            font=\sffamily\bfseries]
        at (current page.center) {\grothWatermark};
    \end{tikzpicture}%
  \fi}

\AtBeginDocument{%
  \begin{center}
    {\large\bfseries Manuscrits de mathématiciens (Gallica) ---
      \ifx\grothShelfmark\empty\grothFolder\else\grothShelfmark\fi}\\[2pt]
    {\itshape\grothFoldertitle}\\[4pt]
    {\small vues \grothFirst--\grothLast
      \ifx\grothBatch\empty\else\ (lot \grothBatch)\fi}\\[2pt]
    \ifx\grothDating\empty\else
      {\small\color{gray!60!black}Datation du catalogue : \grothDating}\\[2pt]
    \fi
    \ifx\grothWatermark\empty\else
      {\small\bfseries\color{red!45!gray}\grothWatermark}\\[2pt]
    \fi
    {\footnotesize\color{gray!60!black}Transcription établie d'après les images IIIF de Gallica.
     Source gallica.bnf.fr / Bibliothèque nationale de France.\\
     Les lectures incertaines sont soulignées, les illisibles notées [\,\ldots\,],
     les ajouts entre crochets ; \textsf{v.} numérote les vues de Gallica,
     \textsf{f.} la foliotation au crayon de la BnF.}
  \end{center}
  \vspace{1em}}

\endinput


\folder{fr-9115}
\shelfmark{Français 9115}
\batch{26}
\pages{501}{520}
\dating{1801-1900}
\watermark{Lecture automatique\\non vérifiée}
\foldertitle{Papiers de Sophie GERMAIN. — Recueil de dissertations et problèmes mathématiques et physiques.}

\begin{document}

\note{Numérotation des feuillets. Chaque page de cette mise au net porte
en tête, au milieu, entre parenthèses, sa pagination : (91) à la vue 501,
puis (92) à (110) aux vues 502 à 520, une par vue, sans lacune et sans
rature. Les feuillets sont écrits des deux côtés, et seul le côté de
pagination impaire porte en haut à droite, à l'encre, d'une écriture plus
ronde, un nombre de la série qui court dans tout le volume : 246 (vue
501), 247 (503), 248 (505), 249 (507), 250 (509), 251 (511), 252 (513),
253 (515), 254 (517), 255 (519). Cette série numérote des feuillets et non
des pages ; elle occupe la place de la foliotation de la Bibliothèque,
mais elle est tracée à la plume et non au crayon, et aucun « f. » n'est
donc porté. Tous ces nombres ont été lus sur la vue ; aucun n'est déduit.
Chaque feuillet est monté sur une feuille de garde, dont le haut ou le bas,
blanc, paraît sur beaucoup de vues. Aucune vue de ce lot n'est blanche ni
répétée.}

\page{501}

et dans celui des $y$, \struck{les sons} on prend pour terme de comparaison
les sons qui accompagnent les figures qui appartiennent a
l'intégrale (C) et \ill{}, \struck{comme je l'ai dit,} \struck{qui} tant d'après
l'opinion de M\textsuperscript{r} Chladni que d'après la théorie que j'ai
cherché a établir, doivent être les mêmes que les sons
donnés par la simple lame dont les extremités sont
libres en vertu des conditions $\frac{ddz}{dx^2}=0$ et $\frac{d^3z}{dx^3}=0$,
on trouve que le son de toutes les figures périodiques
qui renferment des périodes complettes, est \struck{trop} \add{plus} grave
d'environ 1 ton que la théorie ne l'indique.

C'est ainsi, par exemple, que le son de $0|4$ fig. 72--73
est sol$^{\sharp}3$ et celui de $4|3$ fig. 80 est fa$^{\sharp}5$, tandis qu'il
ne devroit y avoir que deux octaves entre ces deux sons,
puisque, d'après la théorie d'Euler, le premier
est proportionel au nombre $\frac{25}{4}$, et que d'après la
théorie des surfaces périodiques, le second est proportionel
au nombre $25$.

Mais \struck{non seulement} le son des figures qui contiennent ou
sont censées contenir une ou plusieurs periodes complettes,
ne sont \add{pas seulement} plus graves d'un ton environ que ne semble l'exiger
la comparaison aux sons donnés par la simple lame;
\struck{mais} cette difference doit \add{encore} être considérée comme un minimum
\struck{du quel} \add{dont} se rapprochent \add{d'autant plus} les sons produits par les vibrations
\note{En tête, au milieu, « (91) », entre parenthèses, non barré ; en haut
à droite, à l'encre, d'une écriture plus ronde, « 246 ». La première
ligne continue une phrase commencée à la vue 500, hors de ce lot. Après
« (C) et », un groupe surchargé et noirci, surmonté d'un point, ne se lit
pas. Dans « le son de toutes », « le » est écrit sur « la ». Les
notations « $0|4$ » et « $4|3$ » sont celles des figures de Chladni ; le
dernier chiffre de « $4|3$ » est tracé comme celui de « 73 ». Les
additions « pas seulement », « encore » et « d'autant plus » sont
écrites au-dessus de la ligne, la première appelée par deux petits
signes en forme de 7. Le « s » et le « t » initiaux se ressemblent dans
cette main ; « son » est lu partout où le sens est celui d'un son
particulier, « ton » là où il s'agit de l'intervalle (« 1 ton », « un ton
environ »). La phrase se poursuit à la vue 502.}

\page{502}

tournantes \struck{se rapprochent d'autant plus qu'elles} \add{ou les figures qui leur appartiennent} contiennent
un plus grand nombre de lignes nodales. C'est ainsi que d'une part
les sons qui accompagnent les fig 63 et 68, \add{bien qu'ils} \struck{ne} devroient
différer que de 2 octaves \struck{different} puisqu'ils sont entr'eux
d'après la théorie $::2:8$, \struck{et qu'ils} different pourtant de 2 octaves
1 ton et $\frac{1}{2}$ ton; desorte que le son dela fig. 63 est
plus grave de 2 tons et $\frac{1}{2}$ ton que ne le veut la
comparaison aux sons delasimple lame, \add{et} que les sons
qui accompagnent les fig. 66\textsuperscript{a} et 69, présentent les mêmes
differences, \struck{et} \add{tandis} que \add{d'autre part} le son de la fig. 74 \struck{et ceux} est
moins eloigné du rapport donné par la théorie, et \struck{qu'enfin}
les sons des fig. 79b et 86 se rapprochent encore plus des
nombres que j'ai cru pouvoir leur assigner en leur appliquant
l'intégrale (I), \struck{d\ill{} même que} \add{car} le son de la fig 86
qui est ut.6 comparé à celui dela fig 75 qui est ut 5
(V table P. 152) est exactement cequ'il doit être pour la comparaison
au cas des périodes complettes, puisque le nombre que
j'ai assigné au premier est 36, et que la théorie donne
18 pour le second.

Des differences du même genre se font remarquer dans la
comparaison des sons des figures qui se forment sur les plaques
rectangulaires differentes du quarré desorte que les figures
qui sur ces diverses plaques presentent une ou plusieurs
périodes complettes donnent toujours, aumoins en prenant au
\note{En tête, au milieu, « (92) », non barré ; aucun nombre en haut à
droite : c'est le dos du feuillet de la vue 501. La première ligne
achève la phrase de la vue 501 (« les sons produits par les vibrations /
tournantes ») ; « se rapprochent d'autant plus qu'elles » y est barré et
remplacé, au-dessus de la ligne, par « ou les figures qui leur
appartiennent ». Le « ne » barré devant « devroient » est surchargé ;
« bien qu'ils » est écrit au-dessus. « $::2:8$ » : le signe entre 2 et
8 a une barre verticale collée aux deux points. « le son de la fig. 74 »
est corrigé de « les sons », l'article surchargé et le « s » final
barré d'une petite croix ; de même « du rapport » corrigé de « des » ;
la fin de « accompagnent » et de « rapprochent » est aussi surchargée.
Le groupe barré après « (I), » commence par « d » et se termine par
« même que » ; le milieu est noirci. « ut.6 » et « ut 5 » : le 5 a la
forme de celui de « fa$^{\sharp}$5 » à la vue 501. La phrase se poursuit
à la vue 503.}

\page{503}

resultat moyen, des sons un peu plus graves qu'il ne le faudroit
pour établir l'accord \struck{ent} dela théorie et de l'experience
Dans la comparaison \uncertain{au} son des figures qui se rapportent aux
mouvemens de la simple lame, et que les figures qui sur les
mêmes plaques ne contiennent point de périodes complettes,
donnent des sons qui dans la même comparaison se trouvent
de beaucoup trop graves.

On pourroit à la rigueur regarder les différences dont
il s'agit ici comme peu importantes, puisqu'il est rare qu'elles
excèdent une tierce, distance qui n'a pas empêché
M\textsuperscript{r} Chladni de \struck{regarder} \add{considérer} comme équivalentes les fig 64 et
65 que la théorie montre aussi devoir appartenir à une
même valeur de l'angle $\underline{\omega}$. Cependant la constance
du phénomène le rend remarquable, et ne permet guères
de le regarder comme produit par les causes physiques
qui influent sur les expériences les mieux faites; car
s'il en étoit ainsi, les differences devroient en quelquesorte
se balancer tandis qu'au contraire elles sont toutes, ou
presque toutes dans le même sens.

Puisqu'il \struck{me paroit} \add{semble} inévitable de conclure \add{de cequi précède} l'existence d'une
cause tendante à rendre les sons qui accompagnent les
figures périodiques, plus graves que la comparaison avec
les sons \add{qui accompagnent} les figures non périodiques \struck{qui peuvent se
former sur les mêmes surfaces}, ne semble l'exiger,
\note{En tête, au milieu, « (93) » ; en haut à droite, à l'encre, d'une
écriture plus ronde, « 247 ». La première ligne continue la phrase de la
vue 502 (« aumoins en prenant un / resultat moyen »). Le mot barré après
« l'accord » est court et commence par « en » ; « Dans », à la ligne
suivante, a une majuscule bien que la phrase continue. Après
« comparaison », « au » est écrit sur un autre mot, peut-être « du ».
« de cequi précède » est écrit au-dessus de la ligne et appelé par un
petit signe sous « conclure ». « tendante » porte à la fin une petite
croix, qui annule semble-t-il un « s » ; « les » devant « figures non
périodiques » est écrit sur « des ». La phrase se poursuit à la vue 504.}

\page{504}

il \struck{\ill{}} \add{me sera permis de} presenter a cet egard une conjecture qui ne me paroît
pas entièrement denuée de probabilité.

On se rappelle que les differences qu'il s'agit d'expliquer
sont à leur \emph{maximum} lorsque la surface périodique
ne presente que des demies periodes tant dans le sens
des $\underline{x}$ que dans celui des $\underline{y}$, c'est ceque j'ai déja
nommé quart de période \struck{pour distinguer l'espace qu'il}
en reservant la denomination \struck{q} de demie période
pour les \struck{figure} \add{parties} composées d'une periode dans le sens d'un
des axes et d'une demie periode seulement dans le
sens de l'autre axe. Cela posé, on se souvient encore
que les differences, en prenant toujours un resultat moyen
entre plusieurs \struck{faits}, sont d'autant moindres que la
surface périodique qui ne comprend pas de periode
complettes est composée d'un plus grand nombre
de demies periodes, et qu'enfin lorsqu'outre les quarts
de périodes et les demies périodes qui se trouvent
dans \struck{toutes} les figures périodiques dont les extremités
sont libres, la figure renferme des periodes complettes
les mêmes differences diminuent encore desorte, \struck{\ill{}} elles
sont \add{même} souvent audessous d'un ton qui peut être considéré
comme leur mesure moyenne.
\note{En tête, au milieu, « (94) », non barré ; aucun nombre en haut à
droite : c'est le dos du feuillet de la vue 503. Au début, « il » est
surchargé et suivi d'un mot court barré ; « me sera permis de » est écrit
au-dessus. « periode » (quart de période, periode complettes) et
« figure » (la figure renferme) portent sur l'« s » final la petite
croix qui annule un pluriel : le singulier est donné. Le « q » isolé
barré après « denomination » est un début de mot abandonné. Dans « pour
les parties composées », « figure » est souligné ou barré d'un trait
léger, et « parties » est écrit au-dessus ; le « les » est surchargé.
« sont » devant « libres » est écrit sur un autre mot. Le groupe barré
après « desorte, » est noirci ; « même » est écrit au-dessus de
« souvent ». Les deux tiers inférieurs de la marge et le bas du feuillet
sont blancs.}

\page{505}

\struck{si on observe que si il s'agissoit de surfaces appuyées de
tout côté les mêmes équations, seroient applicables}

En réfléchissant sur la gradation du phénomène, il paroîtra
peut-être vraisemblable qu'il ait pour cause l'existence de ces
quarts de periode et demi periodes dans les figures qui
se forment sur les plaques dont tous les côtés sont libres.
En effet on voit que si les quarts de periode sont seuls
leur influence doit être a leur maximum qu'elle doit diminuer
par l'addition de demie periodes qui comprennent des periodes
entières dans le sens d'un des axes, et qu'enfin elle doit
se trouver en partie contrebalancée et parconséquent réduite
a son \emph{minimum} lorsque la figure contient des périodes
complettes; desorte que \struck{la théorie ne pourroit s'accorder} dans cette
hypothèse on ne devroit obtenir un accord parfait entre les sons
indiqués par la théorie et les sons donnés par l'experience
que dans le seul cas des surfaces appuyés detout côté.

Malheureusement M\textsuperscript{r} Chladni ne s'est pas occupé de ce
genre de surfaces qui \struck{presentent} \add{ne peuvent être soumises à l'experience qu'avec} beaucoup de difficultés
\struck{a obtenir, et je ne puis appuyer \ill{}}. Le cas \struck{analogue}
des lames libres à leurs extrémités en vertu des conditions
$\frac{ddz}{dx^2}=0$ et $\frac{d^3z}{dx^3}=0$, \struck{si} s'il avoit été essayé avec succès
\uncertain{m'avoit} fourni un fait analogue à celui des surfaces periodiques
libres; et si il eut donné un son plus grave que
le cas correspondant dans la lame dont les extrémités sont
appuyés, j'aurois été autorisé a regarder l'effet des demies
\note{En tête, au milieu, « (95) » ; en haut à droite, à l'encre, d'une
écriture plus ronde, « 248 ». Le feuillet est monté plus bas que les
précédents sur sa feuille de garde, dont le haut, blanc, occupe le
premier quart de la vue. Les deux premières lignes sont barrées d'un
trait continu ; la page commence au paragraphe « En réfléchissant ».
Dans « gradation » et « paroîtra » une lettre est surchargée d'une
tache ; « periode » porte deux fois la petite croix qui annule le « s »
du pluriel, et le singulier est donné. « elle » dans « qu'elle doit
diminuer » est repassé. Le mot barré après « qui », « presentent », est
remplacé par l'ajout au-dessus de la ligne. La fin de la ligne barrée
suivante, après « appuyer », ne se lit pas (deux mots courts, le second
terminé en « -dit »). Dans $\frac{ddz}{dx^2}$, le premier $d$ et
l'exposant sont noircis d'une tache. Le premier mot de « m'avoit
fourni » peut se lire « n'avoit » ; le sens demande « m'avoit » ou
« m'auroit ». La phrase se poursuit à la vue 506.}

\page{506}

periodes comme constaté. Mais l'experience n'ayant rien décidé à
cet egard, je suis réduit à présenter l'explication que je viens
de donner comme une simple conjecture.

En effet quoique les cas dela corde tendue et de la flute
ouverte aux deux bouts semblent présenter deux faits
comparables et susceptibles, parconséquent, de fournir aumoins
des inductions, \struck{dans} sur l'objet présent, ils ne peuvent
pourtant être d'aucun secours, \struck{parceque} \add{car} la détermination
du son dans les instrumens à vent ne peut encore être
regardée comme parfaitement rigoureuse.

Voila, ce me semble, ici la difficulté; L'ingénieuse hypothèse
que M\textsuperscript{r} De la Place a imaginé et que M\textsuperscript{r} Poisson
a ensuite introduite dans le calcul, suffit à la verité
pour expliquer comment la vitesse réelle du son surpasse
celle que l'analyse indique, parceque cette vitesse qui n'a
guere été mesurée que sur des sons graves, est regardée
\struck{p} tacitement par les physiciens comme independante de
la place \struck{qu'occupe} \add{que} le son transmis occupe dans
l'échelle diatonique; desorte qu'en considérant d'une part
cette \struck{tons} \add{vitesse} comme uniforme et de l'autre l'influence
du dégagement de la chaleur comme egalement uniforme
on peut satisfaire aux conditions dela question.

Mais s'il s'agissoit de décider dans la correction à
appliquer aux sons d'une flûte, \struck{l'effet la quantité}
\note{En tête, au milieu, « (96) », non barré ; aucun nombre en haut à
droite : c'est le dos du feuillet de la vue 505. La première ligne
achève la phrase de la vue 505 (« l'effet des demies / periodes comme
constaté »). Un petit trait oblique, au-dessus de « à présenter », est
peut-être un signe d'insertion sans ajout. « dans », devant « sur l'objet
présent », est chargé de traits plus épais et paraît annulé. Dans « du
son dans les instrumens », l'« s » final porte la petite croix qui
annule un pluriel. Devant « tacitement », un « p » isolé est barré.
« cette » est corrigé de « ces » par une tache, et « tons » qui suivait
est barré, « vitesse » écrit au-dessus. La phrase se poursuit à la vue
507.}

\page{507}

l'influence du coëfficient \struck{de} \add{proportionnel a} la chaleur degagée, et celle de l'existence
des demies périodes, on seroit d'autant plus embarassé que
l'on ne \struck{pourroit regarder} peut affirmer la constance d'un
coëfficient qui est fondé sur un fait variable dans son
amplitude et dans sa frequence, \struck{dont mesuré pendant un
tems donné desorte que l'on auroit} et que parconséquent
deux effets dont on n'auroit pas de mesure exacte, se
trouveroient confondus.

Quoiqu'il en soit du degré de confiance que merite
l'explication que je viens de soumettre au jugement dela
classe, je n'ai pas cru devoir dissimuler la difficulté
qui \struck{est} en est l'objet. Mais il me sera permis d'observer
qu'avant l'explication proposée par M\textsuperscript{r} De la Place,
la théorie des instrumens àvent et celle de la transmission
du son, qui \struck{présentoient} donnoient des resultats fort
éloignés de ceux de l'experience, n'en étoient pas moins
\struck{admis} \add{regardées} par les géomètres et les physiciens comme entièrement
satisfesantes, \struck{et} que \add{d'ailleurs} la difference que j'ai essayé d'expliquer
est bien peu considérable puisqu'elle \add{dans son \emph{maximum}} excède \struck{\ill{} fort} \add{à} peine
\struck{peu} celles \struck{que} \add{dont} M\textsuperscript{r} Chladni n'a pas tenu compte dans
la formation de sa table p. 152, et qu'elle n'est vraiment
remarquable que par sa constance.

Je vais m'occuper àprésent des plaques rectangulaires differentes du
quarré.
\note{En tête, au milieu, « (97) » ; en haut à droite, à l'encre, d'une
écriture plus ronde, « 249 ». La première ligne achève la phrase de la
vue 506 (« dans la correction à appliquer aux sons d'une flûte, /
l'influence du coëfficient… »). « proportionnel a » est écrit au-dessus
de « dela », dont le « de » est barré. « fondé » et « variable » portent
à la fin la petite croix qui annule une lettre (« fondée »,
« variables ») ; la forme corrigée est donnée. Le début d'« amplitude »
est surchargé. Dans « la difference », l'article est repassé sous
« d'ailleurs », écrit au-dessus. Après « excède », un mot court surchargé
puis « fort » sont barrés ; un accent seul, au-dessus, donne « à » ;
« dans son maximum » est écrit au-dessus de la ligne, « maximum »
souligné. « celles » s'accorde avec un pluriel qui n'est pas écrit. Le
texte s'arrête au tiers inférieur de la page sur « du quarré. » ; le
reste du feuillet est blanc.}

\page{508}

§.5 \emph{Comparaison dela théorie à l'experience dans le cas des plaques
rectangulaires differentes du quarré}

N.\textsuperscript{o} 17. L'intégrale (C) N.\textsuperscript{o} 3, dont j'ai longuement developpé l'application
aux plaques quarrées, (V. N.\textsuperscript{os} 10 11 et 12 §.4) peut aussi servir à rendre
compte d'une partie des figures qui se \struck{mon} montrent sur les plaques
vibrantes \struck{differentes du quarre} rectangulaires differentes du quarré, \struck{aussi}
\struck{bien que} \add{et en même tems aussi} des sons qui accompagnent chacune des mêmes figures.

Cependant, en fesant attention à la nature de la \struck{formule (C)} \add{question},
on verra aisément qu'elle exige que la valeur de l'angle $\underline{\omega}$
soit lamême pour chacune des fonctions de $\underline{x}$ et $\underline{y}$ qui la
\struck{composent} \add{entrent dans la formule (C)}. Ainsi, \struck{si} dans la plus-part des cas, on aura necessairement
$B=0$ \add{ou $\mathcal{A}=0$}, c'est adire qu'il ne pourra se former \struck{que des lignes
alors} sur les plaques rectangulaires dont les mouvemens sont
compris dans l'intégrale (C), que de simples lignes droites
nodales parallelles entr'elles; \struck{et} On ne \struck{pourra obtenir} \add{trouvera}
des \struck{va} figures susceptibles de \emph{distorsions} que pour les
valeurs de \add{l'angle} $\underline{\omega}$ qui \struck{s'accorderont} auront entr'elles le rapport
necessaire pour compenser la difference entre la grandeur
des côtés des plaques \struck{d} que l'on considère.

Je vais éclaircir ceci par des exemples. Si on a une plaque
dont les côtés soient entr'eux $::5:7$. on verra aisément
que le nombre \struck{qui convient au cas de \emph{3} lignes parallelles
étant 25} auquel le son doit être proportionnés dans le cas
de 3 lignes droites parallelles, etant $\frac{25}{4}$, et \struck{celui auquel}
le nombre auquel le \struck{pr} son doit être proportionnel.
\note{En tête, au milieu, « (98) », non barré ; aucun nombre en haut à
droite : c'est le dos du feuillet de la vue 507. Le titre est souligné
et précédé de « §.5 » ; le signe § est tracé comme un S suivi d'un
point. Dans « (V. N.\textsuperscript{os} 10 11 et 12 §.4) », le « 12 »
est à demi effacé. « rapport » et « necessaire » portent à la fin la
petite croix qui annule le « s » du pluriel ; le singulier est donné.
Le $B$ de « $B=0$ » est écrit sur une autre lettre, noircie, et le signe
$=$ est repassé ; l'ajout « ou $\mathcal{A}=0$ » est au-dessus de la
ligne, la lettre étant la capitale bouclée que le lot 22 a rendue par
$\mathcal{A}$. Le dénominateur de $\frac{25}{4}$ est surchargé et
peut se lire 8. « proportionnés » est gardé tel qu'il est écrit. Le
texte s'arrête à mi-page sur « doit être proportionnel. », phrase
inachevée ; le bas du feuillet est blanc. La suite, si elle existe, est
à la vue 509.}

\page{509}

dans le cas de quatre lignes nodales parallelles, étant $\frac{49}{4}$,
3 lignes nodales parallelles entr'elles et transversales au côté qui est
comme 7, doit donner le même son que 4 lignes nodales parallelles
entr'elles et transversales au côté qui est comme 5. Il suit dela
que la valeur de l'angle $\underline{\omega}$ est lamême pour l'un et
l'autre de ces cas, et qu'il peut exister alors differentes
figures \struck{sur qui seront toutes accompagnées du même son} \add{équivalentes}
puisque la supposition $B=0$ cesse d'être necessaire.

Independamment de la considération de l'intégrale (C), M\textsuperscript{r}
Chladni paroît avoir senti que ce ne pouvoit être que dans
le cas dela compensation dont j'ai parlé plus haut, qu'il
pouvoit exister sur les plaques rectangulaires differentes du
quarré, des figures nodales \struck{differentes} \add{diverses} entr'elles, mais toutes
equivalentes à un certain nombre de points de repos dans
la simple lame dont les extremités sont libres. \struck{(V. T. d'acoustique
N.\textsuperscript{o} 136).} \struck{et} On peut voir N.\textsuperscript{o} 136 du traité d'acoustique, que
dans le cas que j'ai pris pour exemple, parcequ'il repond
a la première des figures observées sur les plaques rectangulaires
l'auteur regarde en effet les figures \struck{qui} \struck{a} qui se sont montrées,
comme \emph{intermédiaires} entre les deux especes de vibrations, qu'il atribue
à la plaque partagée par quatre lignes nodales \add{parallelles} suivant une
de ses dimensions, et par 3 pareilles lignes suivant l'autre
de ses dimensions. N.\textsuperscript{o} 121 du même traité, on peut encore voir que
la plaque dont les côtés sont entr'eux $::3:5$, présente aussi
des figures intermediaires aux deux manières de vibrer qui repondent
aux suppositions successives $\mathcal{A}=0$ et $B=0$, dans l'intégrale
(C). car alors, suivant la remarque de l'auteur,
\note{En tête, au milieu, « (99) » ; en haut à droite, à l'encre, d'une
écriture plus ronde, « 250 ». Comme à la vue 505, le feuillet est monté
bas sur sa feuille de garde, dont le haut blanc occupe le premier quart
de la vue. La première ligne continue la phrase laissée en suspens au
bas de la vue 508 (« le nombre auquel le son doit être proportionnel /
dans le cas de quatre lignes nodales parallelles ») : le point qui la
fermait là n'arrêtait donc pas la phrase. Le dénominateur de
$\frac{49}{4}$ est surchargé. « équivalentes » est écrit au-dessus de
« sur », en tête de la ligne barrée ; « puisque » est écrit sur un
premier « puis ». « paroît » et « pouvoit » (« ce ne pouvoit être »)
sont surchargés. Le renvoi barré au « T. d'acoustique » et le renvoi
gardé portent le même numéro, dont le chiffre du milieu est le 3 à tête
plate de cette main. Dans « les figures », l'« s » de l'article et
celui du nom sont surchargés ; « qui » et « a » qui suivent sont barrés.
« N.\textsuperscript{o} 121 » est repassé. Dans « $::3:5$ », le
second chiffre a la forme du 5 de la vue 501. $\mathcal{A}$ est la
capitale bouclée de la vue 508. « atribue » est son orthographe. La
phrase se poursuit à la vue 510.}

\page{510}

2 lignes nodales parallelles dans un des sens doit donner le même
son que 4 pareilles lignes dans l'autre sens.

Demême encore \struck{pour} sur la plaque dont les côtés sont entr'eux
$::3:1$. Les suppositions successives $\mathcal{A}=0$ et $B=0$ qui repondent
\struck{et}, en suivant la notation de M\textsuperscript{r} Chladni, aux cas $5|0$ et $0|2$
donnent les deux \struck{mêmes} figures auxquelles les fig. 77 sont
equivalentes (V. T. d'acoustique N 122).

\marginal{Parmi les fig. pl. 6 du traité \add{d'acoustique}, celles qui repondent
aux trois cas dont j'ai parlé, et les differentes
fig. 170 sont les seules
qui presentent plusieurs
formes dependantes
\struck{de \ill{}} toutes
de l'intégrale (C).}
\struck{Parmi les figures qui ont été \add{pl. 6 du traité d'acoustique} observées sur les plaques rectangulaires
celles qui appartiennent aux trois cas dont je viens de parler
et les fig 170 sont les seules qui puissent être expliquées a l'aide de
l'intégrale (C) et} Rien ne seroit plus aisé que de
suivre de point en point, pour chacune d'elles la
formation des differentes lignes nodales qui les composent.
mais j'aurois à répéter des détails fastidieux qui peut être
\uncertain{ont} déja été que trop rebattus dans la description des
figures que lamême intégrale (C) donne sur la plaque
quarrée. ainsi je me contenterai de presenter a cet

\struck{Cependant, afin de montrer comment la supposition $\mathcal{A}=B$,
appliquée aux plaques rectangulaires dont les côtés sont
inégaux entr'eux peut donner lieu a des figures
dependantes du rapport entre les mêmes côtés, \add{figures très} et très
differentes parconséquent de celles qui ont été expliquées
pour la plaque quarrée; je crois devoir entrer dans
le détail de la formation dela fig. 163b}
\note{En tête, au milieu, « (100) », non barré ; aucun nombre en haut à
droite : c'est le dos du feuillet de la vue 509. Le feuillet est monté
bas sur sa feuille de garde. La première ligne achève la phrase de la
vue 509 (« suivant la remarque de l'auteur, / 2 lignes nodales… ») ; le
premier chiffre peut se lire 3, le second, surchargé, est lu 4 ; la fin
de « parallelles » et de « doit » est surchargée. Dans « $5|0$ », le 5 a
la forme de celui de la vue 501. La note de la marge de gauche, d'une
plume plus fine, remplace le passage barré qui lui fait face (« Parmi
les figures… (C) et ») ; « d'acoustique » y est écrit au-dessus de
« traité », dont l'« s » final est barré d'une petite croix, et le
groupe barré avant « toutes » commence par « de » et ne se lit pas
au-delà. Dans le passage barré, « pl. 6 du traité d'acoustique » est
écrit au-dessus de « ont été ». « ainsi je me contenterai de presenter a
cet » reste inachevé au bout de la ligne ; un petit signe en forme de
virgule est sous « quarrée ». « ont » (« qui peut être ont déja été »)
est écrit sur un autre mot. Tout le dernier paragraphe, de
« Cependant » à « fig. 163b », est barré ligne par ligne ; « figures
très » est écrit au-dessus de « et très », lui-même barré. Le bas du
feuillet est blanc.}

\page{511}

egard quelques reflexions générales et de les appliquer,
par forme d'exemple, à la figure 163\emph{b} que je supposerai
tracée.

Lorsque les suppositions successives $\mathcal{A}=0$ et $B=0$ donnent
lieu, l'une et l'autre, à \struck{des} \add{la formation de} nombres pairs de lignes
nodales droites parallelles entr'elles, les valeurs de $f(x)$ et de
$f(y)$ \struck{considérées} prises à l'origine \struck{a} \struck{sont} \add{peuvent être} egales entr'elles
parceque toutes deux sont dela forme $2(1+e^{\omega})$; ces deux
valeurs \struck{sont} \add{peuvent} encore \add{être} egales entr'elles lorsque les suppositions
$\mathcal{A}=0$ \struck{et} $B=0$ donnent lieu l'une et l'autre àlaformation
de \struck{ligne} nombres impairs delignes nodales droites; car alors
ces deux fonctions sont dela forme $2(1-e^{\omega})$; Mais si les
mêmes suppositions, savoir $\mathcal{A}=0$ et $B=0$, donnent lieu, l'une
à la formation d'un nombre impair de lignes nodales
parallelles entr'elles, et que l'autre entraîne laformation
d'un nombre pair de pareilles lignes, on ne pourra
plus regarder $f(x)$ et $f(y)$ comme égales entr'elles à
l'origine, \struck{car} quoique la valeur de $\underline{\omega}$ soit egale
de part et d'autre.

Il résulte de ceque je viens de dire, que la supposition
$\mathcal{A}=-B$ qui dans aucun cas, \add{excepté celui du quarré,} ne peut donner lieu
àla formation delignes nodales parallelles \uncertain{et} diagonales
\struck{ne peut même ne peut même donner lieu a la formation}
\note{En tête, au milieu, « (101) » ; en haut à droite, à l'encre, d'une
écriture plus ronde, « 251 ». La première ligne achève la phrase
inachevée de la vue 510 (« de presenter a cet / egard »). Le numéro de
la figure est « 163 » suivi d'un signe souligné qui peut se lire « b »
ou « 6 » ; c'est la figure dont la vue 510 annonçait, dans le passage
barré, « le détail de la formation ». « la formation de » est écrit
au-dessus de « des », lui-même surchargé d'une petite croix. Le mot entre
« parallelles » et « diagonales » est surchargé et peut-être barré. Le
bas de la page porte une ligne entièrement barrée, et la phrase reste
en suspens ; elle se poursuit à la vue 512.}

\page{512}

sur les plaques rectangulaires, ne peut même expliquer laformation
d'une ligne nodale ayant son origine au point $\underline{a}$ que
dans le seul cas où les suppositions $\mathcal{A}=0$ et $B=0$ entraînent
l'une \struck{et} \struck{l'ont} l'autre la formation de nombres pairs
ou celle de nombres impaires de lignes nodales droites.

La fig. 175\emph{b} de M\textsuperscript{r} Chladni présente un exemple
de ce cas; \struck{par} et les suppositions $\mathcal{A}=0$ et $B=0$
donnent \add{en effet} $2|0$ et $0|4$.

Si dans le cas où les suppositions successives
$\mathcal{A}=0$ et $B=0$ donnent lieu, l'une à un
nombre pair, l'autre à un nombre impair
de lignes nodales droites, on vouloit avoir
une figure dont une des lignes nodales commençat
au point $\underline{a}$ d'origine, aulieu de la supposition
$\mathcal{A}=-B$, il faudroit faire celle-ci
$\mathcal{A}(1-e^{\omega}) = -(1+e^{\omega})B$. La fig. 163\emph{b} (V. la table)
paroît appartenir à cette \struck{supposition} \add{hypothèse}; \struck{et} en la supposant
tracée, on voit qu'elle satisfait à toutes les conditions
que la théorie exige. En effet à partir du point
$\underline{a}$ la ligne nodale s'avance diagonalement vers le point
$i$ qui lui est opposé parceque \struck{f(x) et f(y) qui}
\struck{etoient} \add{les quantités} $(1-e^{\omega})f(x)$ et $(1+e^{\omega})f(y)$ qui étoient egales au
point $\underline{a}$, diminuent en même tems jusqu'à ceque elles \struck{de}
\note{En tête, au milieu, « (102) », non barré ; aucun nombre en haut à
droite : c'est le dos du feuillet de la vue 511. La première ligne
continue la phrase de la vue 511 (« qui dans aucun cas […] ne peut
donner lieu àla formation delignes nodales parallelles et diagonales /
sur les plaques rectangulaires, ne peut même expliquer… »). Une petite
croix au-dessus de « ou », en tête de la ligne « ou celle de nombres
impaires », semble signaler un ajout ; elle n'appelle aucune note.
« fig. 175\emph{b} » et « fig. 163\emph{b} » : le dernier
signe, souligné (rendu en \emph{}), peut se lire « b » ou « 6 ». « $2|0$ et $0|4$ » : la
notation des figures de Chladni. Au-dessus de « à toutes », un mot
ajouté et barré, peut-être « aussi ». Le dernier mot de la page,
« de », est barré ; la phrase se poursuit à la vue 513. Le bas du
feuillet, sous cette ligne, est coupé par le bord de la feuille de
garde : rien n'y est écrit.}

\page{513}

deviennent nulles \struck{en même tems} \add{àlafois} au point $i$. À partir de ce point
la ligne entre dans un intervalle où \add{les plus grandes valeurs} des mêmes quantités
$(1-e^{\omega})f(x)$ et $(1+e^{\omega})f(y)$ sont égales \struck{entr'elles}, ou cequi revient aumême
quant a l'effet physique, presqu'égales entr'elles. et c'est pourcela
que, suivant ceque j'ai expliqué N.\textsuperscript{o} 11, àl'occasion
dela quatrième valeur del'angle $\underline{\omega}$, la figure nodale
montre quatre branches autour d'un point $\underline{m}$ de cet
intervalle. Il en est demême pour les autres intervalles
où se trouvent de semblables points $\underline{m}$; et quant au
changement de direction dela courbe, j'ai eu
occasion tant de fois d'en expliquer la cause que
je me contenterai de renvoyer \struck{au même N.\textsuperscript{o} 11}
\struck{dans lequel} au passage du N.\textsuperscript{o} 11 que j'ai cité plus
haut. \struck{et} Il me suffira \add{donc} de remarquer que conformément
à la théorie, la presente figure nodale passe par
les 12 points de \emph{repos invariables}.

Il est évident que les autres figures 163 données par
M\textsuperscript{r} Chladni, remplissent également cette condition;
car l'une d'elles montre quatre lignes nodales sinueuses
ayant chacune une \emph{flexion} et $\frac{1}{2}$, et l'autre
presente 3 lignes nodales sinueuses ayant chacune
2 \emph{flexions} complettes; et on se rappelle qu'une flexion
\struck{1} $\frac{1}{2}$ passe par trois points de repos invariables,
tandis que deux \emph{flexions} complettes passent par
\note{En tête, au milieu, « (103) » ; en haut à droite, à l'encre, d'une
écriture plus ronde, « 252 ». Le feuillet est monté bas sur sa feuille
de garde. La première ligne achève la phrase de la vue 512 (« jusqu'à
ceque elles / deviennent nulles »). La fin de « deviennent » est
surchargée. « intervalle » et « physique » portent sur l'« s » final la
petite croix qui annule un pluriel ; le singulier est donné. « les plus
grandes valeurs » est écrit au-dessus de la ligne, et « des » est écrit
sur « les » ; « entr'elles », ajouté au-dessus de « égales », est barré.
La fin d'« autour » est noircie. « Il », après « haut. », est écrit sur
un « et » barré. Dans « une flexion $\frac{1}{2}$ », le signe barré avant
la fraction paraît être un 1. « repos invariables » est souligné, comme
« flexion » et « flexions ». La phrase se poursuit à la vue 514.}

\page{514}

quatre pareils points de repos invariables.

On peut établir comme une regle générale pour juger
aquelle valeur de l'angle $\underline{\omega}$ se rapporte une figure
nodale composée de lignes sinueuses, que si
$\underline{m}$ est le nombre des lignes droites nodales auxquelles
la supposition $\mathcal{A}=0$ donne lieu, et que $\underline{n}$
soit le nombre des lignes droites nodales qu'entraîne
la supposition $B=0$, la figure sera
necessairement composée de $\underline{m}$ lignes sinueuses
ayant chacune $\frac{n}{2}$ \emph{flexions}, ou de $\underline{n}$ lignes
également sinueuses dirigées dans le sens opposé
aux premiers, et presentant chacune $\frac{m}{2}$ flexions.

On voit donc que la considération des points de
repos invariables àlaquelle j'ai déja eu recours plusieurs
fois, est du plus grand secours pour l'explication
des figures qui sont assujetties à passer par
chacun d'eux; mais je me dispenserai defaire
de nouvelles applications de la regle précédente, et
je passerai toutdesuite àla considération des intégrales
periodiques.

N.\textsuperscript{o} 18 \emph{surfaces periodiques}. Pour appliquer les intégrales
données N.\textsuperscript{os} 4 et 5 au cas des surfaces rectangulaires
differentes du quarré, il faut multiplier
\note{En tête, au milieu, « (104) », non barré ; aucun nombre en haut à
droite : c'est le dos du feuillet de la vue 513. Le feuillet est monté
bas sur sa feuille de garde. La première ligne achève la phrase de la
vue 513 (« deux flexions complettes passent par / quatre pareils
points »). Le $B$ de la seconde supposition est écrit sur une autre
lettre, peut-être la capitale bouclée $\mathcal{A}$. « flexions » est
souligné la première fois, non la seconde ; « surfaces periodiques »,
titre du N.\textsuperscript{o} 18, est souligné. « aux premiers » est
écrit ainsi pour « aux premières ». La fin de « toutdesuite » est
surchargée. La phrase se poursuit à la vue 515.}

\page{515}

les ordonnées parallelles au plus petit côté par la fraction renversée
qui exprime le rapport entre la longueur des deux côtés dela
plaque rectangulaire, desorte que si, par exemple, les longueurs
\struck{des} \add{deux} côtés sont entr'elles $::s:t$ $t$ etant $<$ que $s$, la forme
generale de ces equations sera
\[ z = \ldots\ldots\,\text{sin}.\Big(\frac{\pi.tmx}{tA}\Big)\,\text{sin}\Big(\frac{\pi.sny}{tA}\Big). \]
Cette formule montre qu'il y aura $\underline{m}$ lignes nodales parallelles
à l'axe des $\underline{y}$ et $\underline{n}$ lignes nodales parallelles à l'axe des
$\underline{x}$, et que, si on changeoit $m$ en $n$, la figure seroit
alors fort differente dela premiere.

Elle montre aussi que le son doit en général être
proportionnel au nombre $\frac{t^2m^2+s^2n^2}{t^2}$.

M\textsuperscript{r} Chladni n'a pas jugé necessaire defaire connoître
\struck{la} \struck{p} la totalité des figures qu'il a observées sur les plaques
rectangulaires differentes du quarré; mais celles qu'il a
données presentent un caractere important pour essayer
vérifier la théorie, et je regrette de n'avoir pas assez
de tems pour entrer dans le détail de leur description
individuelle. Cependant, en me bornant même à des
généralités, j'espère pouvoir rendre un compte suffisant
des singularités que presentent les divers groupes
d'équations que M\textsuperscript{r} Chladni regarde comme
équivalents.
\note{En tête, au milieu, « (105) » ; en haut à droite, à l'encre, d'une
écriture plus ronde, « 253 ». Le feuillet est monté sur sa feuille de
garde, dont une bande blanche occupe le haut de la vue. La première
ligne achève la phrase de la vue 514 (« il faut multiplier / les
ordonnées… »). « deux » est écrit au-dessus de « des », dont la fin est
surchargée. Dans la formule, la lettre qui précède $ny$ au second
numérateur est écrite sur une autre et se lit $s$ ; les deux
dénominateurs sont $tA$, et les points de suspension sont les siens,
comme au lot 22 (« = … », facteur laissé en blanc). La formule est
gardée telle qu'elle est écrite. Dans $\frac{t^2m^2+s^2n^2}{t^2}$, les
exposants sont écrits au-dessus des lettres, d'une plume plus fine.
Devant « totalité », « la » est surchargé et un « p » est barré.
« équivalents » est écrit sur « équivalentes ». Le bas du feuillet est
blanc.}

\page{516}

Jusqu'ici je n'ai employé, pour \struck{rendre compte} \add{expliquer} des distorsions
des figures nodales que la considération d'une seule forme
d'intégrale, desorte que ç'a toujours été par des changemens
\struck{exécutés} effectués dans une seule et même équation, que
j'ai crû pouvoir expliquer le changement des lignes nodales
lorsque le son restoit le même.

Mais par la nature du genre de calcul qui convient
àlaquestion que je traite, non seulement les differentes
valeurs de $\underline{z}$ qui résultent des diverses intégrales
particulieres, satisfont à l'équation differentielle, mais
la somme de plusieurs de ces valeurs y satisfait
également; et rien n'empêche de supposer en effet
que $\underline{z}$ soit ainsi composé. Cependant, comme il \struck{faut}
ne s'agit que des mouvemens qui doivent produire
un son mesurable, les valeurs de $\underline{z}$ composées comme
je viens de le dire, ne peuvent être employées
que dans les cas où le son des differens termes
produits par l'addition des diverses valeurs de
$\underline{z}$ qui viennent des intégrales particulieres
composantes, est le même pour chacune d'elles.

Pour effectuer l'application des formules composées
d'après le principe que je viens d'établir
\note{En tête, au milieu, « (106) », suivi d'un point, non barré ; aucun
nombre en haut à droite : c'est le dos du feuillet de la vue 515. Le
feuillet est monté sur sa feuille de garde, dont le haut et le bas
paraissent. La page commence un paragraphe nouveau. Après « employé »,
une virgule ou un « s » final. « des » (« des distorsions ») est écrit
sur « les ». « le son » est corrigé de « les sons » : l'« s » de
l'article et celui du nom portent la petite croix qui annule le
pluriel. « composé » (« que z soit ainsi composé ») : le « z » est
surchargé. La phrase reste en suspens au bas du texte, qui s'arrête aux
deux tiers de la page ; elle se poursuit à la vue 517.}

\page{517}

il suffira de multiplier chacun des termes par \struck{une} une constante
arbitraire, et de tracer les figures qui conviendront à la supposition
de toutes ces constantes \add{nulles}, excepté une \struck{supposée nulle}. La rencontre
des differentes figures produites par les suppositions
successives des divers coëfficiens nuls, donnera un certain
nombre de points \struck{de} invariables de repos, qui,
comme je l'ai dit ailleurs, devront appartenir
à toutes les figures renfermées dans la valeur composée
de $\underline{z}$. \struck{mais} \add{et} Non seulement il \struck{peut} existera
\add{$\times$ comme dans les cas qui ont été expliqués jusqu'ici,} des points invariables
de repos, $\times$ mais on \uncertain{pourra} trouver \struck{aussi} \add{alors} de veritables lignes de repos invariables.
\marginal{par exemple \struck{et} si pour plus de simplicité}
\struck{Si} on se borne au cas où la valeur de $\underline{z}$ ne seroit
\struck{si pour plus de simplicité,}
formée que dela somme de deux \struck{valeurs} \add{termes} données
par deux intégrales particulieres differentes, on sentira
aisément qu'il pourra arriver \add{en effet} que les deux figures
appartenantes à chacune de ces deux valeurs ayent
une ligne nodale commune, et qu'alors \struck{la} cette
ligne sera une veritable ligne invariable de
repos, desorte qu'elle sera commune àtoutes
les distorsions des figures appartenantes à
la valeur de $\underline{z}$ ainsi composée.

A l'egard des conditions des limites, on voit ici
une nouvelle preuve de la necessité de distinguer
l'état analytique des extrémités physiquement
\note{En tête, au milieu, « (107) », suivi d'un point ; en haut à
droite, à l'encre, d'une écriture plus ronde, « 254 ». Le feuillet est
monté bas sur sa feuille de garde, dont le haut blanc occupe le premier
quart de la vue. La première ligne achève la phrase de la vue 516
(« d'après le principe que je viens d'établir / il suffira… »). Le
premier « une » devant « constante » est noirci. « nulles » est écrit
au-dessus de « excepté une », et « supposée nulle » qui suivait est
barré. Les deux croix sont les siennes : celle qui est écrite
au-dessus de la ligne, devant « comme dans les cas qui ont été expliqués
jusqu'ici », appelle cet ajout, écrit d'une plume plus fine ; celle qui
suit « de repos, » en marque la place. Le verbe devant « trouver » est
surchargé et lu avec doute ; « alors » est écrit au-dessus de « aussi »,
barré. La note de la marge de gauche, « par exemple si pour plus de
simplicité », remplace le « Si » barré en tête de la ligne et le « si
pour plus de simplicité, » barré sous la ligne ; un mot court y est
barré après « exemple », peut-être « et ». Dans « on se borne »,
l'« s » final de « bornes » porte la petite croix qui l'annule. Le
texte s'arrête au bas de la page au milieu d'une phrase ; elle se
poursuit à la vue 518.}

\page{518}

libres; car rien n'empêche de composer la valeur de $\underline{z}$
de termes qui satisferont \struck{aux} aux conditions des extrémités
en vertu de differens couples d'équations, desorte
que le mot \emph{libre} n'aura a proprement parler
ni l'une ni l'autre des significations que je lui
ai reconnues jusqu'ici, mais que sa valeur sera
composée des deux premieres

Pour appliquer la théorie que je viens d'établir aux
cas observés par M\textsuperscript{r} Chladni, je reprendrai la
formule $\frac{t^2m^2+s^2n^2}{t^2}$ àlaquelle j'ai trouvé que le
son qui \struck{accompagnent} accompagne les figures \struck{qui se forment}
\struck{nodal}des plaques rectangulaires, devoit être proportionnel
en fesant $s=6$ $t=5$, $m=4$ et $n=1$ cette formule
devient $\frac{16.25+96}{25} = \frac{4(109)}{25}$.

en fesant toujours $s=6$ et $t=5$, mais prenant $m=2$, $n=3$,
on a par la même formule $\frac{4.25+9.36}{25} = \frac{4(106)}{25}$.
L'intervalle entre les sons qui appartiennent à ces deux
cas, est donc $\frac{109}{106}$ c'est àdire beaucoup moindre qu'un
demi ton mineur et guere plus grand qu'un \emph{comma}.
La limite des erreurs possibles dans les experiences les
mieux faites, et parconséquent aussi l'influence des
causes physiques inappréciables, est donc plus
grande que la distance entre les sons des
\note{En tête, au milieu, « (108) », non barré ; aucun nombre en haut à
droite : c'est le dos du feuillet de la vue 517. Le feuillet est monté
bas sur sa feuille de garde, dont le haut blanc occupe le premier
quart de la vue. La première ligne achève la phrase de la vue 517
(« des extrémités physiquement / libres »). Le premier « aux » est
écrit sur un autre mot et répété. « premieres » n'a pas de point. Dans
« accompagnent », la finale « -ent » est barrée ; un trait vertical
sépare « les » de « figures ». En tête de la ligne suivante, « des » est
écrit sur un mot commencé, qui paraît être « nodal- » : la lecture en
est incertaine. Le terme « 96 » du premier numérateur : le premier
chiffre a la boucle fermée d'un 9 ; le calcul ($s^2n^2 = 36$, et
$16\cdot25+36 = 436 = 4\cdot109$) demande 36, le 3 de cette main ayant
une queue qui le fait prendre pour un 9 ; la page est gardée telle
qu'elle se lit. Le second calcul, $4\cdot25+9\cdot36 = 424 = 4\cdot106$,
est juste. Le « 1 » de « 109 » au membre de droite est surchargé.
« comma » est souligné. Une petite croix à l'encre pâle, sous
« appartiennent », n'appelle rien. La phrase se poursuit à la vue 519.}

\page{519}

deux \struck{cas} dont j'ai parlé, et ces cas peuvent parconséquent
produire des figures \struck{comp} résultantes de la
somme des termes qui composent les valeurs de
$\underline{z}$, convenables a chacun d'eux.
\marginal{D'ailleurs \struck{\uncertain{Au reste}} le son de \uncertain{$4|1$}
qui appartient au cas des
vibrations tournantes, peut, comme
on \uncertain{sait}, être plus grave que ne
l'indique le \struck{\ill{}} nombre que donne
la théorie, et la difference entre
109 et 106 est dans ce sens puis
que le dernier nombre est $<$ que le 1\textsuperscript{er}}
C'est ceque M\textsuperscript{r} Chladni a observé (V. N.\textsuperscript{o} 114 fig. 157)
et parceque l'une et l'autre valeur de $\underline{n}$ donne
une ligne droite nodale parallelle au plus grand
côté dela plaque et passant par le centre de
la figure, cette ligne sera une veritable ligne
invariable de repos et devra se trouver, comme
elle se trouve en effet dans toutes les figures
qui appartiennent à la même forme de \add{l'angle} $\underline{\omega}$.

Dans le cas précédent les deux termes de
$\underline{z}$ étoient donnés par une même equation
périodique mais ils peuvent être produits également
, comme je l'ai déjà dit, par des intégrales
entièrement differentes entr'elles

C'est ainsi par exemple que suivant la remarque de
M\textsuperscript{r} Chladni (V. N.\textsuperscript{o} 115) lorsque les côtés dela plaque
sont entr'eux $::5:4$ les sons \add{qui accompagnent} \struck{des} figures $5|0$ et
$1|4$ sont les mêmes et \add{que ces fig.} \struck{presentent} \add{montrent} des \struck{differences} \add{variations} qu\ill{} \struck{\uncertain{cit}} l'auteur
\note{En tête, au milieu, « (109) » ; en haut à droite, à l'encre, d'une
écriture plus ronde, « 255 ». La première ligne achève la phrase de la
vue 518 (« la distance entre les sons des / deux cas dont j'ai
parlé »). « cas », en tête, est surchargé et semble barré. La note de
la marge de droite, d'une plume plus fine, est placée en face de « C'est
ceque M\textsuperscript{r} Chladni a observé » et des lignes qui
suivent ; aucun signe ne l'appelle dans le texte. Son premier mot,
barré, est lu avec doute « Au reste » ; « D'ailleurs » est écrit
au-dessus. La figure de Chladni qu'elle nomme a pour second chiffre un
trait vertical qui peut être un 1 ou un 3 mal formé ; la vue 501 citait
« $4|3$ fig. 80 ». Le mot après « comme on » est surchargé. Le mot
barré devant « nombre » est court et ne se lit pas. À la ligne de
« valeur de $\underline{n}$ », une grosse tache d'encre précède « de ».
« l'angle » est écrit au-dessus de $\underline{\omega}$, lui-même écrit
sur une autre lettre. La virgule qui ouvre la ligne « , comme je l'ai
déjà dit » est bien en tête de ligne. À la dernière ligne, « qui
accompagnent » est écrit au-dessus de « des », barré ; « que ces fig. »
au-dessus de « montrent », lui-même écrit au-dessus de « presentent »,
barré ; « variations » au-dessus de « differences », barré. La fin de la
ligne est surchargée : après « qu », une lettre ou un accent illisible,
puis un groupe barré qui paraît être « cit », puis « l'auteur ». La
phrase se poursuit à la vue 520.}

\page{520}

qu'il a représentées par les figures 158. \struck{On voit en}
effet en \struck{repassant} mettant $s=5$ $t=4$, $m=1$ et $n=4$ dans
la formule $\frac{t^2m^2+s^2n^2}{t^2}$, elle donne $\frac{\ill{}.16+25.16}{16} = \struck{\frac{\ill{}}{16}} = 26$
pour le nombre auquel le son de $1|4$ doit être proportionnel
\struck{et que la théorie d'Euler} \struck{fo} d'après la théorie d'Euler,
le nombre auquel $5|0$ est proportionnel \struck{étoit}
$\frac{81}{4}.\frac{25}{16}$; \struck{d'où il suit} \add{et parconséquent} que les sons \struck{\ill{} doivent égaux}
qui appartiennent à l'un et l'autre cas, doivent être
\struck{égaux entr'eux, car leur rapport est exprimé par
la fraction 1125} entr'eux $::2025:1664$. c'est-àdire
à peu près $::5:4$. \struck{q} Le rapport précédent indique une
différence de 2 tons environ entre les sons que l'on
compare; et en se rappellant que les vibrations
tournantes \struck{presentent} \add{sont} toutes plus graves \struck{de \ill{}}
d'un intervalle qui a pour limite d'une part 1 ton
et de l'autre 2 tons et $\frac{1}{2}$ ton, que ne semble
l'exiger leur comparaison avec les sons dela simple lame
on sentira que la réunion des termes qui conviennent
a $5|0$ et a $1|4$ doit donner lieu à des figures
compliquées par leur influence simultanée

Si le terme dû à $5|0$ donnoit, \struck{s} s'il étoit seul,
5 lignes nodales droites \struck{et} parallelles entr'elles \struck{et}
\note{En tête, au milieu, « (110) », non barré ; aucun nombre en haut à
droite : c'est le dos du feuillet de la vue 519. Le feuillet est monté
sur sa feuille de garde, dont le bas blanc occupe le dernier tiers de
la vue. La première ligne achève la phrase de la vue 519 (« des
variations qu… l'auteur / qu'il a représentées par les figures 158 ») ;
« qu'il » y est repassé. Dans « $m=1$ » et « $n=4$ », les chiffres sont
écrits sur d'autres, noircis. Le premier facteur du numérateur, devant
« .16 », est surchargé et ne se lit pas (le calcul, $t^2m^2 = 16$,
demande 1) ; la fraction intermédiaire, barrée, a un numérateur de trois
chiffres surchargés (on lit « 281 », là où $16+400 = 416$). Le « 81 »
de $\frac{81}{4}$ est repassé. Le mot barré devant « doivent égaux » est
court et ne se lit pas ; le groupe barré après « plus graves »
commence par « de ». « 2025 : 1664 » est juste
($\frac{81}{4}\cdot\frac{25}{16} = \frac{2025}{64}$ et $26 =
\frac{1664}{64}$) ; la fraction barrée « 1125 » n'a pas de dénominateur.
« compliquées » est écrit sur « compliqués », « simultanée » sur une
autre fin de mot. La page s'arrête sur « entr'elles », suivi d'un « et »
barré ; la phrase se poursuit au-delà de ce lot, à la vue 521.}

\end{document}
